% WDDA Vorlesungsnotizen Preamble
% Geteiltes Styling für alle Notizen-Dokumente

% ---- Seitenlayout ----
\geometry{margin=2.5cm}

% ---- Typografie ----
\usepackage{lmodern}
\usepackage{microtype}

% ---- Code-/Verbatim-Zeilenumbruch ----
% Lange Zeilen in Code-Chunks am Rand umbrechen
\usepackage{fvextra}
\DefineVerbatimEnvironment{Highlighting}{Verbatim}{commandchars=\\\{\},breaklines,breakanywhere}
\fvset{breaklines=true,breakanywhere=true}

% ---- Mathematik ----
\usepackage{amsmath, amssymb}

% ---- Mehr Abstand um Grafiken ----
% Grössere Float-Abstände (für Bilder mit Caption innerhalb von figure-Floats)
\setlength{\textfloatsep}{16pt plus 2pt minus 2pt}
\setlength{\intextsep}{14pt plus 2pt minus 2pt}
\setlength{\floatsep}{14pt plus 2pt minus 2pt}
% Auch um inline-includegraphics (knitr-Plots ohne fig.cap) etwas Luft schaffen
\usepackage{etoolbox}
\AtBeginDocument{%
  \let\WDDAoldincludegraphics\includegraphics
  \renewcommand{\includegraphics}[2][]{%
    \par\addvspace{0.8em}%
    \WDDAoldincludegraphics[#1]{#2}%
    \par\addvspace{0.8em}%
  }%
}

% ---- Mehr Abstand um Code-Blöcke ----
% Pandoc verpackt knitr-R-Chunks in 'Shaded' (Highlighting) und R-Output in 'verbatim'
\BeforeBeginEnvironment{Shaded}{\par\addvspace{0.8em}}
\AfterEndEnvironment{Shaded}{\par\addvspace{0.8em}}
\BeforeBeginEnvironment{verbatim}{\par\addvspace{0.6em}}
\AfterEndEnvironment{verbatim}{\par\addvspace{0.6em}}

% ---- Farben (passend zum Foliendeck) ----
\usepackage{xcolor}
\definecolor{BFHdunkelblau}{RGB}{0,60,120}
\definecolor{BFHorange}{RGB}{237,128,0}
\definecolor{BFHmittelblau}{RGB}{0,90,170}
\definecolor{BFHmittelgruen}{RGB}{0,130,80}

% ---- Hyperlinks ----
\usepackage{hyperref}
\hypersetup{
  colorlinks=true,
  linkcolor=BFHdunkelblau,
  urlcolor=BFHmittelblau,
  citecolor=BFHmittelgruen
}

% ---- Eigene Befehle ----
\newcommand{\textti}[1]{\texttt{#1}}
\newcommand{\courseName}{WDDA}
\DeclareMathOperator{\Var}{Var}
\DeclareMathOperator{\SD}{SD}
\DeclareMathOperator{\SE}{SE}
\DeclareMathOperator{\Cov}{Cov}
\DeclareMathOperator{\corr}{corr}

% ---- Boxen für Kernkonzepte ----
\usepackage{tcolorbox}
\newtcolorbox{keybox}[1][]{
  colback=BFHdunkelblau!5,
  colframe=BFHdunkelblau,
  fonttitle=\bfseries,
  title=#1
}
\newtcolorbox{examplebox}[1][]{
  colback=BFHorange!5,
  colframe=BFHorange,
  fonttitle=\bfseries,
  title=#1
}

% ---- Section-Styling ----
\usepackage{titlesec}
\titleformat{\section}{\Large\bfseries\color{BFHdunkelblau}}{\thesection}{1em}{}
\titleformat{\subsection}{\large\bfseries\color{BFHmittelblau}}{\thesubsection}{1em}{}

% ---- Kopf-/Fusszeile ----
\usepackage{fancyhdr}
\pagestyle{fancy}
\fancyhf{}
% Kurzform der Section im Header: ohne Nummer, nicht in Grossbuchstaben
\renewcommand{\sectionmark}[1]{\markboth{#1}{}}
\fancyhead[L]{\small\textcolor{BFHdunkelblau}{WDDA}}
\fancyhead[R]{\small\textcolor{BFHdunkelblau}{\nouppercase{\leftmark}}}
\fancyfoot[C]{\thepage}
\renewcommand{\headrulewidth}{0.4pt}
